\documentclass[8pt]{extarticle}
\usepackage[margin=0.55in]{geometry}
\usepackage{booktabs}
\usepackage{array}
\usepackage{times}
\usepackage[T1]{fontenc}
\usepackage{caption}
\usepackage{amssymb}
\usepackage{graphicx}

\setlength{\parindent}{0pt}
\captionsetup{font=small,skip=2pt,singlelinecheck=false}

\makeatletter
\newcommand{\rows}[1]{\@@input /Users/vasyl/zadumai/legalbench_map/results_cgd/#1_rows.tex }
\makeatother
\newcommand{\header}{%
\textbf{Task} & \textbf{LLM} & \textbf{`Free' (CPU run) NLP alone} & \textbf{LLM alone} &
\textbf{Confidence cascade (same budget)} & \textbf{CGD} & \textbf{$\Delta$ vs better} &
\textbf{LLM-call share} & \textbf{Outcome} \\}

\begin{document}

\begin{center}
{\Large\bfseries Class-Gated Deferral (CGD): Spending LLM Calls Only Where the LLM Is Better}\\[0.3em]
{\normalsize Vasyl Rakivnenko, \ldots}\\[0.15em]
{\small September 2026}
\end{center}
\vspace{0.35em}

\begin{center}
\begin{minipage}{0.74\linewidth}
\begin{center}{\bfseries Abstract}\end{center}
\noindent
We introduce class-gated deferral (CGD), to our knowledge the first cascade
algorithm that decides class by class, which model should answer: a
computationally `free' (CPU-run) NLP classifier
trained on the task's labels, or a large language model (LLM). For every
class the classifier can predict, CGD checks which of the two models is more
accurate on held-out data and hands that whole class to the winner, so LLM
calls are spent only where the LLM is actually better. Uncertainty-based
cascades, such as Online Cascade Learning~[5], work differently: they send
the LLM whatever the free model is least sure about, even when the LLM is
worse on those items. On
LegalBench~[1], across nine task--LLM combinations, CGD beats both the free
classifier and the LLM in seven, and it beats a confidence cascade with the
same LLM budget in six, by up to 17 points of balanced accuracy. On
CUAD~[3] contract clauses with a production API LLM, CGD again beats both
models and outperforms a same-budget cascade by 4.8--8.1 points. On
Financial PhraseBank~[4], it keeps the majority neutral class on the free
classifier, cuts LLM calls by 60\%, and beats both components in 16 of 28
LLM setups. The gain
grows with how class-dependent the LLM's competence is ($r = 0.57$ across
59 pairs), which suggests a simple test for when to gate by class.
\end{minipage}
\end{center}

\vspace{1.0em}

\begin{table}[h]
\caption{LegalBench~[1], balanced accuracy (LegalBench's metric).
\emph{Confidence cascade} gives the LLM the same share of items as CGD,
escalating the free classifier's least-confident ones.
opp115\_data\_retention comes from the OPP-115 privacy-policy corpus~[2].}
\label{tab:legal}
\centering
\resizebox{\linewidth}{!}{%
\begin{tabular}{@{}l l r r r r r r c@{}}
\toprule
\header
\midrule
\rows{table1_legal}
\bottomrule
\end{tabular}%
}
\end{table}

\vspace{1.4em}
\begin{table}[h]
\caption{CUAD~[3] contract clauses with a production API LLM. Accuracy
(the test sets are class-balanced, so this equals balanced accuracy).}
\label{tab:cuad}
\centering
\resizebox{\linewidth}{!}{%
\begin{tabular}{@{}l l r r r r r r c@{}}
\toprule
\header
\midrule
\rows{table2_cuad}
\bottomrule
\end{tabular}%
}
\end{table}

\vspace{1.4em}
\begin{table}[h]
\caption{Financial PhraseBank~[4], accuracy. CGD keeps the neutral class on
the free classifier and hands positive and negative items to the LLM
(40\% of calls). Six of 28 LLM setups shown; CGD beats both components in
16 of the 28.}
\label{tab:fpb}
\centering
\resizebox{\linewidth}{!}{%
\begin{tabular}{@{}l l r r r r r r c@{}}
\toprule
\header
\midrule
\rows{table3_fpb}
\bottomrule
\end{tabular}%
}
\end{table}

\vspace{1.0em}
\footnotesize
\noindent \textbf{`Free' (CPU run) NLP alone}: a classifier trained on each
task's labeled data (5-fold cross-validation for LegalBench and CUAD; a
separate training split for Financial PhraseBank) that runs on CPU.
\textbf{CGD} decides which predicted classes go to the LLM from both models'
accuracy on the training folds; all routing decisions are out-of-fold,
averaged over 10 fold splits ($\times$3 cross-validation repeats for
LegalBench and CUAD).
$^{*}$\,$p<0.05$ vs.\ the better single model; $^{\dagger}$\,$p<0.05$ vs.\ the
same-budget cascade (paired McNemar test for accuracy, paired bootstrap for
balanced accuracy). \emph{Tie}: within 0.5 points.

\vspace{1em}
\normalsize
\noindent{\bfseries References}\\[0.3em]
\noindent[1] N.~Guha, J.~Nyarko, D.~E.~Ho, C.~R\'e, et al. ``LegalBench:
A Collaboratively Built Benchmark for Measuring Legal Reasoning in Large
Language Models.'' arXiv:2308.11462, 2023.\\[0.3em]
[2] S.~Wilson, F.~Schaub, A.~A.~Dara, F.~Liu, S.~Cherivirala, P.~G.~Leon,
et al. ``The Creation and Analysis of a Website Privacy Policy Corpus.''
ACL 2016, pp.~1330--1340.\\[0.3em]
[3] D.~Hendrycks, C.~Burns, A.~Chen, S.~Ball. ``CUAD: An Expert-Annotated
NLP Dataset for Legal Contract Review.'' NeurIPS 2021 Datasets and
Benchmarks Track.\\[0.3em]
[4] P.~Malo, A.~Sinha, P.~Korhonen, J.~Wallenius, P.~Takala. ``Good Debt or
Bad Debt: Detecting Semantic Orientations in Economic Texts.'' Journal of
the Association for Information Science and Technology 65(4):782--796,
2014.\\[0.3em]
[5] L.~Nie, Z.~Ding, E.~Hu, C.~Jermaine, S.~Chaudhuri. ``Online Cascade
Learning for Efficient Inference over Streams.'' ICML 2024, PMLR
235:38071--38090.

\end{document}
